\documentclass[11pt,a4paper]{article}
\usepackage[T2A]{fontenc}
\usepackage[utf8]{inputenc}
\usepackage[russian]{babel}
\usepackage{geometry}
\usepackage{amsmath,amssymb,mathtools}
\usepackage{array,booktabs,longtable}
\usepackage[dvipsnames]{xcolor}
\usepackage{tikz}
\usepackage[colorlinks=true,linkcolor=blue!55!black]{hyperref}
\geometry{left=19mm,right=19mm,top=20mm,bottom=21mm}
\setlength{\parindent}{0pt}
\setlength{\parskip}{4pt}
\setlength{\emergencystretch}{2em}
\renewcommand{\arraystretch}{1.22}
\newcommand{\reviewpart}[1]{\par\smallskip\textbf{#1}\par\smallskip}

\begin{document}
\begin{titlepage}
  \thispagestyle{empty}
  \begin{center}
    \vspace*{0.28\textheight}
    {\LARGE\bfseries Ревью домашней работы\par}
    \vspace{15mm}
    {\large Автор: Лёвин Артём Александрович\par}
    \vspace{8mm}
    {\large 24 октября 2026 года\par}
  \end{center}
  \vfill
  \begin{center}
    \begin{tikzpicture}
      \draw[blue!35!black,line width=0.65pt]
      (0,0) .. controls (1.9,0.42) and (3.4,-0.42) .. (5.3,0)
            .. controls (7.1,0.42) and (8.6,-0.42) .. (10.4,0);
    \end{tikzpicture}
  \end{center}
\end{titlepage}

\clearpage
\section*{Сводная таблица ревью}
\small
\setlength{\tabcolsep}{3pt}
\begin{longtable}{@{}>{\raggedright\arraybackslash}p{0.9cm}>{\raggedright\arraybackslash}p{3.2cm}>{\raggedright\arraybackslash}p{5.0cm}>{\raggedright\arraybackslash}p{2.7cm}>{\raggedright\arraybackslash}p{2.9cm}@{}}
\toprule
\textbf{№} & \textbf{Тема} & \textbf{Суть недочёта} & \textbf{Ваш ответ} & \textbf{Правильный ответ}\\
\midrule
\endhead
\hyperref[task:6]{6} & Деление десятичной дроби на обыкновенную & Не объяснено, почему при записи дроби $0{,}08$ появляется знаменатель $100$. & $\frac15$; пояснение отсутствует & $\frac15$; пояснение о сотых долях\\
\addlinespace[4pt]
\hyperref[task:10]{10} & Линейное уравнение & Некорректная промежуточная запись; не выполнена требуемая проверка подстановкой. & $10\frac25$; без проверки & $10\frac25$; с проверкой\\
\bottomrule
\end{longtable}
\normalsize

\clearpage
\thispagestyle{plain}
\vspace*{\fill}
\begin{center}
  {\LARGE\bfseries Разбор ошибок}
\end{center}
\vspace*{\fill}
\clearpage

\phantomsection\label{task:6}
\section*{Задание № 6}
\reviewpart{Текст задания}
Вычислите $0{,}08:\dfrac25$. Объясните, почему знаменатель при переводе десятичной дроби равен $100$.

\reviewpart{Ваше решение}
В записи выполнен переход от десятичной дроби к обыкновенной, затем произведено деление:
\[
0{,}08:\frac25=\frac8{100}:\frac25=\ldots=\frac15.
\]
Значение частного найдено правильно, однако требуемого словесного объяснения знаменателя нет.

\reviewpart{Ваш ответ}
$\dfrac15$. Обязательное пояснение не приведено.

\reviewpart{\textcolor{red}{Ошибка:} пропущен обязательный элемент ответа}
Недочёт состоит не в конечном числовом результате, а в отсутствии объяснения, почему у дроби $\dfrac8{100}$ знаменатель именно $100$.

\reviewpart{Где именно возникает недочёт}
Последний корректный существенный шаг:
\[
0{,}08=\frac8{100}.
\]
Сразу после него решение переходит к вычислению частного. Именно здесь необходимо было письменно объяснить происхождение знаменателя $100$. Такого объяснения в работе нет.

\reviewpart{Почему решение неполное}
Цифры после запятой показывают, на какие равные доли разделена единица. У числа $0{,}08$ \textbf{две цифры после запятой}. Первая позиция отвечает за десятые доли, а вторая~--- за сотые. Поэтому $0{,}08$ означает восемь сотых, то есть $\dfrac8{100}$.

Число $100$ не выбирается произвольно: каждая из двух позиций после запятой даёт множитель $10$ в знаменателе. Таким образом, $10\cdot10=100$. Без этого обоснования не выполнено прямое требование условия, даже если вычисление правильное.

\reviewpart{\textcolor{blue!60!black}{Ключевая идея:} десятичные разряды и деление дробей}
Две цифры после запятой означают сотые доли единицы. Для деления на обыкновенную дробь следует умножить делимое на дробь, обратную делителю.

\reviewpart{Исправление от последнего правильного шага}
Сохраняем верно записанное равенство $0{,}08=\dfrac8{100}$ и добавляем недостающее объяснение: \emph{после запятой стоят две цифры, следовательно, число состоит из сотых долей единицы; поэтому знаменатель равен ста}. Далее можно продолжить вычисление.

\reviewpart{Полное правильное решение}
\textbf{Шаг 1.} Переведём десятичную дробь в обыкновенную. Число $0{,}08$~--- это восемь сотых, потому что после запятой две цифры:
\[
0{,}08=\frac8{100}=\frac2{25}.
\]
При сокращении числитель и знаменатель разделены на $4$.

\textbf{Шаг 2.} Заменим деление на $\dfrac25$ умножением на обратную дробь $\dfrac52$:
\[
\frac2{25}:\frac25=\frac2{25}\cdot\frac52.
\]

\textbf{Шаг 3.} Перемножим числители и знаменатели и сократим результат:
\[
\frac2{25}\cdot\frac52=\frac{2\cdot5}{25\cdot2}=\frac{10}{50}=\frac15.
\]

\reviewpart{Проверка}
Умножим полученное частное на делитель. Должно получиться делимое:
\[
\frac15\cdot\frac25=\frac2{25}=\frac8{100}=0{,}08.
\]
Проверка подтверждает численный результат.

\reviewpart{\textcolor{teal!60!black}{Итог:}}
\[
\boxed{0{,}08:\frac25=\frac15.}
\]
Знаменатель $100$ возникает потому, что запись $0{,}08$ содержит две цифры после запятой и обозначает восемь сотых.

\reviewpart{Задача на закрепление}
Вычислите $0{,}06:\dfrac35$. Объясните, почему при переводе десятичной дроби $0{,}06$ в обыкновенную знаменатель равен $100$.

\clearpage
\phantomsection\label{task:10}
\section*{Задание № 10}
\reviewpart{Текст задания}
Решите уравнение $\dfrac14x-2=\dfrac35$. Выполните проверку подстановкой.

\reviewpart{Ваше решение}
В работе после исходного уравнения присутствует промежуточная строка вида
\[
\frac14x-2+\frac35=2+\frac35.
\]
Затем записаны преобразования
\[
\frac14x=\frac{13}{5},\qquad
x=\frac{13}{5}:\frac14=\frac{52}{5}=10\frac25.
\]
Конечное значение $x$ правильное. Однако первая промежуточная строка не является равносильным переходом от исходного уравнения, а требуемая проверка подстановкой отсутствует.

\reviewpart{Ваш ответ}
$\displaystyle x=10\frac25$. Проверка не приведена.

\reviewpart{\textcolor{red}{Ошибка:} неверно записан перенос слагаемого и пропущена проверка}
Число $-2$ нужно устранить из левой части \emph{прибавлением $2$ к обеим частям} исходного уравнения. В промежуточной записи вместо этого слева сохранено $-2$ и добавлено $\dfrac35$, тогда как справа стоит $2+\dfrac35$. Такое равенство нельзя получить из исходного одинаковым действием с двумя частями.

Отдельный недочёт~--- отсутствие обязательной проверки найденного корня подстановкой.

\reviewpart{Где именно возникает ошибка}
Последний безусловно правильный шаг~--- исходное уравнение:
\[
\frac14x-2=\frac35.
\]
Первый некорректный переход~--- строка с $\frac14x-2+\frac35$ в левой части. Она не отражает прибавление одного и того же числа к двум сторонам исходного равенства.

Правильная следующая строка должна быть такой:
\[
\frac14x-2+2=\frac35+2.
\]
Здесь к обеим частям прибавлено именно $2$.

\reviewpart{Почему этот переход неверен}
Равенство сохраняется, когда к его левой и правой частям прибавляют \emph{одно и то же число}. Если прибавить разные числа или изменить выражение только с одной стороны, получится другое уравнение, которое необязательно имеет тот же корень.

В исходном уравнении вычитается $2$. Для устранения этого слагаемого необходимо прибавить $2$, потому что $-2+2=0$. Приписывать $\dfrac35$ в левую часть для такого переноса не нужно. Правильно найденное позднее значение $\dfrac{13}{5}$ не делает некорректную промежуточную строку верной.

\clearpage
\reviewpart{\textcolor{blue!60!black}{Ключевая идея:} равносильные преобразования}
При решении уравнения выполняют одинаковые действия с обеими частями, после чего упрощают выражения. Найденное значение неизвестной подставляют в \emph{исходное} уравнение, а не только в последнюю преобразованную строку.

\reviewpart{Исправление от последнего правильного шага}
Возвращаемся к исходному равенству и прибавляем $2$ к обеим частям:
\[
\frac14x-2+2=\frac35+2
\quad\Longrightarrow\quad
\frac14x=\frac{13}{5}.
\]
Дальнейший способ нахождения $x$ в работе выбран верно. После получения ответа остаётся выполнить подстановку.

\reviewpart{Полное правильное решение}
\textbf{Шаг 1.} Запишем исходное уравнение:
\[
\frac14x-2=\frac35.
\]

\textbf{Шаг 2.} Прибавим $2$ к обеим частям:
\[
\frac14x-2+2=\frac35+2.
\]
В левой части $-2+2=0$, поэтому
\[
\frac14x=\frac35+\frac{10}{5}=\frac{13}{5}.
\]

\textbf{Шаг 3.} Чтобы найти $x$, разделим обе части на $\dfrac14$. Деление на четверть равносильно умножению на $4$:
\[
x=\frac{13}{5}:\frac14
  =\frac{13}{5}\cdot\frac41
  =\frac{52}{5}
  =10\frac25.
\]

\reviewpart{Проверка подстановкой}
Подставим $x=\dfrac{52}{5}$ в \emph{левую часть исходного уравнения}:
\[
\begin{aligned}
\frac14\cdot\frac{52}{5}-2
 &=\frac{52}{20}-2\\
 &=\frac{13}{5}-\frac{10}{5}\\
 &=\frac35.
\end{aligned}
\]
Получилось ровно $\dfrac35$, то есть правая часть исходного уравнения. Следовательно, найденный корень удовлетворяет уравнению.

\reviewpart{\textcolor{teal!60!black}{Итог:}}
\[
\boxed{x=\frac{52}{5}=10\frac25.}
\]
Ответ численно верен. Решение необходимо оформить с правильным равносильным переходом и закончить проверкой подстановкой.

\reviewpart{Задача на закрепление}
Решите уравнение $\dfrac13x-4=\dfrac23$. Выполните проверку подстановкой.

\vfill
\begin{center}
\footnotesize Лёвин Артём Александрович эксклюзивно для Ксении Васильченко
\end{center}
\end{document}
